\begin{remark}[Period variation and norm comparisons]\hypertarget{tag-01Q5}{}\label{fam:src:A.4.3}\label{fam:st:A.4.3}
\StatementRecord{fam:src:A.4.3}{proved}{01Q5}{\SeriesRef{fam:src:A.4.2}, \SeriesRef{fam:src:A.1.1}, \SeriesRef{fam:src:A.2.1}, \SeriesRef{fam:src:A.3.1}; \cite[Proposition 4.7]{ViaclovskyLog2026}; \cite[Lemmas 6.6--6.9 and Proposition 6.15]{Yang2026}}

\begin{enumerate}[label=\arabic*.,leftmargin=*]
\item On a fixed ordinary fibre, the three last periods remain fixed and span $\mathbb C^2$, whereas the first changes by $\delta e$. A marked complex-linear identification fixing those three periods must be the identity, so the marked period map is nonconstant. The set of integral markings of one torus is countable. A holomorphic nonconstant map from a disc has uncountable image; hence the unmarked fibre family is locally non-isotrivial. \ref{fam:src:A.1.1} transfers this to total spaces. This is the period argument of \cite[Proposition~4.7, pp.~14--15]{ViaclovskyLog2026}. A nonconstant classifying map gives positive dimension of its local target. It does not by itself determine the derivative at the base point, as the explicit holomorphic map $z\mapsto z^2$ has nonconstant image and derivative zero at zero. \ref{fam:src:A.4.2} computes the derivative here.

\item The comparison in \cite[Lemmas~6.6--6.7 and~6.9, pp.~31--32; Remark~6.8, p.~32; Proposition~6.15, pp.~37--38]{Yang2026} uses an integral commuting lattice map, an affine fibre lift, normalization of the root base change, and a norm-lattice defect. His Proposition 6.15 obtains $6/r$ from the first-column scaling and its order-six norm image. Here the first-column scaling is one and the two computed norm images impose parity, yielding two. These matrix substitutions identify the commuting maps, finite norm equations and collar corrections with the ones just displayed. The nonzero absolute Kodaira--Spencer conclusion in \cite[Theorem~6.18, p.~39]{Yang2026} is obtained for the present family by the explicit computation and intrinsic-pencil transfer in \cref{fam:src:A.4.2}.
\end{enumerate}
\end{remark}